% ==============================================================================
% pretextual/abstract.tex
% Resumo em língua estrangeira (Abstract)
% ==============================================================================

\begin{resumo}[Abstract]
\begin{otherlanguage*}{english}
This technical report presents the architecture, engineering specification, complete implementation, and automated validation of a hybrid data acquisition, analytics, and actuation pipeline for industrial automation and the Internet of Things (IoT), developed within an IoT Learning Situation (\textit{Situação de Aprendizagem}) at Centroweg Senai. Originally conceived as a Network Time Protocol (NTP) server based on an ESP32 microcontroller and a Global Navigation Satellite System (GNSS) receiver, the project underwent a complete scope pivot after indoor satellite reception constraints and laboratory infrastructure limitations rendered the initial prototype unfeasible, leading to a formally approved transition toward a self-contained and reproducible Industrial IoT (IIoT) testbed. At the 24\,VDC PNP field layer, documented through KiCad~8 and CircuiTikZ schematics, four pushbuttons (\texttt{B1} to \texttt{B4}) simulate part presence, critical level, overtemperature, and normally closed fail-safe emergency stop sensors, while four Light-Emitting Diodes (\texttt{L1} to \texttt{L4}) with 2.2\,k$\Omega$ current-limiting resistors represent the motor drive, relief valve, alarm indicator, and normal operation status. At the control layer, modularly programmed in Structured Text (IEC~61131-3) within CODESYS~V3.5 (\texttt{GVL\_Modbus}, \texttt{FB\_\allowbreak Debounce\allowbreak Counter}, \texttt{FB\_\allowbreak Watchdog\allowbreak Heartbeat}, \texttt{FB\_\allowbreak Safety\allowbreak Interlock}, and \texttt{PLC\_PRG} running at a 10\,ms scan cycle and exported as PLCopen TC6 XML v2.01), the WEG PLC500 Programmable Logic Controller operates as a Modbus TCP/IP server on port 502 with 50\,ms input debouncing, event counters, a heartbeat communication watchdog, and local deterministic safety interlocking. At the orchestration and analytics layer, a five-tab Node-RED~v4 application uses the \texttt{node-red-\allowbreak contrib-\allowbreak modbus} package to poll 23 zero-based registers every 500\,ms (2\,Hz), persist telemetry and audit trails in PostgreSQL~16 (four tables and three operational views), and compute statistical features over 60\,s sliding windows with a 15\,s stride. These feature vectors feed a FastAPI and \textit{scikit-learn} machine learning microservice---comprising an \texttt{Isolation\allowbreak Forest} anomaly detector and a supervised \texttt{HistGradient\allowbreak Boosting\allowbreak Classifier} trained on 18\,h of 2\,Hz telemetry (129,600 raw samples and 3,688 windows separated by $\ge 60$\,s purge gaps)---coupled with a Large Language Model (LLM) diagnostic agent constrained by JSON Schema Draft 2020-12 validation and a monotonicity safety guard. Finally, the complete architecture is verified through automated Software-in-the-Loop testing across ten test cases (\texttt{CT-01} to \texttt{CT-10}) driven by a 10\,ms Modbus TCP digital twin (\texttt{plc\_\allowbreak simulator.py}).

\vspace{\onelineskip}
\noindent
\textbf{Keywords}: Industrial Internet of Things. WEG PLC500. Modbus TCP/IP. Node-RED. Machine learning. Digital twin.
\end{otherlanguage*}
\end{resumo}
